% !PRINTED 657
\weberpage{657}

In this case $\xi$ now runs, by \S.\,172, III, through a complete
system of residues not divisible by $q$ with respect to the modulus
$m_1$, and in (10) we may understand by $\xi'$ the least positive
solution of the congruence
\[
\xi\xi' \equiv 1\ (\mathrm{mod}\ m_1)
\]
because, if we alter the exponents $\xi'$ in (10) with respect to the
modulus $m_1$, only $m$th powers of the prime factors $\mathfrak{p}$
are added.

We can then apply \gesperrt{Kummer's} theorem [\S.\,174, (16)] to
the product (10), replacing $m$ there by $m_1$, and, if we understand
by $r_1 = r^{m_2}$ an $m_1$th root of unity, and by the $\eta$ % [source sic: German prints a full stop after ``weil''; the translation supplies the grammatically required comma]
certain periods of the $p$th roots of unity, we obtain:
\[
\prod^{\xi} \mathfrak{p}_\xi^{\xi'} = (r_1, \eta)^{m_1},
\]
and consequently
\weq{11}{(\prod^{\xi} \mathfrak{p}_\xi^{\xi'})^{hm_2} = (r_1, \eta)^{hm}.}
Accordingly we have the theorem:

6. \gesperrt{If $p$ is a prime and $p-1$ is divisible by $q$, or,
when $q=2$, by 4, then the positive exponent $h$ can be so
determined that the number contained in $\Omega_m$}
\[
\omega (r_1, \eta)^{-hm}
\]
\gesperrt{contains the prime factors of $p$ only to such powers
whose exponent is divisible by $m$.}

It is now also no longer necessary to distinguish Theorem 5. from
Theorem 6.; for when $m_1 = 2$, 6. passes over into 5., because we
then have $r_1 = -1$, and by Bd. I, \S.\,171
\[
(r_1, \eta)^m = (-1, \eta)^m = \sqrt{p}^{\,m}
\]
holds $^1$).

If we combine the result of Theorems 3. to 6. into one formula, we
obtain, if $\varepsilon$ denotes a unit functional, $\varphi$ any
functional of the field $\Omega_m$, and $p$ runs through a sequence
of primes that are congruent to 1 with respect to the modulus $q$,
among which the same prime $p$ may also occur several times,
\weq{12}{\omega = (r, x_0)^m = \varepsilon \varphi^m \prod^{p} (r^{m_2}, \eta)^m.}

\noindent
$^1$) In the author's memoir in Bd. 8 of the Acta mathematica it was
neglected to give special emphasis to the case of Theorem 5.
